\documentclass[11pt]{article}
\usepackage[margin=1in]{geometry}
\usepackage{amsmath,amssymb}
\usepackage{enumitem}
\usepackage{hyperref}
\usepackage{microtype}
\usepackage[T1]{fontenc}
\usepackage[utf8]{inputenc}

\newcommand{\R}{\mathbb R}
\newcommand{\F}{\mathcal F_\lambda}
\newcommand{\Sig}{\Sigma_\lambda}
\newcommand{\dd}{\,d}

\begin{document}

\section{The First Pressure: Small Energy Can Still Look Dangerous}

Fix a possible terminal point \((x_0,T)\). For \(r>0\), write
\[
B_r(x_0)=\{x\in\R^3: |x-x_0|<r\}
\]
for the spatial ball of radius \(r\) around \(x_0\).

The number \(r\) is the spatial scale being tested. Large \(r\) means a coarse neighborhood of \(x_0\); small \(r\) means a closer test of \(x_0\). The limit \(r\to0\) tests whether the apparent danger persists all the way to the point.

The matching time scale is \(r^2\). It comes from the exact Navier-Stokes scaling. Set
\[
x=x_0+ry,
\qquad
t=T+r^2s,
\qquad
u_r(y,s)=r\,u(x,t),
\qquad
p_r(y,s)=r^2p(x,t).
\]
With this change of variables, the Navier-Stokes equation keeps the same form. One unit of rescaled space corresponds to physical length \(r\), and one unit of rescaled time corresponds to physical time \(r^2\). The physical region selected by that scale is
\[
Q_r=B_r(x_0)\times(T-r^2,T).
\]
Equivalently, \(Q_r\) records the flow in \(B_r(x_0)\) over the matching backward time interval before \(T\).

Here \(u(x,t)\) is the velocity field and \(p(x,t)\) is the pressure. At the final time, the ordinary energy seen in the spatial ball is
\[
E(r)=\int_{B_r(x_0)} |u(x,T)|^2\dd x .
\]
This number can go to zero simply because the ball is shrinking. A small ball can contain little total energy while the velocity inside it still looks large at that scale.

The CKN quantity measures that scaled danger on \(Q_r\):
\[
C(r)=r^{-2}\int_{T-r^2}^{T}\int_{B_r(x_0)}
\left(|u(x,t)|^3+|p(x,t)|^{3/2}\right)\dd x\dd t .
\]
The factor \(r^{-2}\) puts different radii in the same scale-critical units. That is why a small amount of raw energy can still matter: the question is how the flow looks after the neighborhood has been zoomed back to unit size.

If \(U_r\) and \(P_r\) denote representative velocity and pressure sizes on
\[
Q_r=B_r(x_0)\times(T-r^2,T)
\]
then the spacetime volume contributes \(r^5\). At the level of scale counting,
\[
\int_{Q_r}\left(|u|^3+|p|^{3/2}\right)\dd x\dd t
\hbox{ has size about } r^5\left(U_r^3+P_r^{3/2}\right),
\]
so the CKN quantity has size about
\[
C(r)\hbox{ has size about } r^3\left(U_r^3+P_r^{3/2}\right).
\]
For \(C(r)\) to stay order-one as \(r\to0\), the representative velocity size is allowed to be order \(r^{-1}\), and the representative pressure size is allowed to be order \(r^{-2}\).

So the raw energy can disappear in the shrinking volume while the rescaled velocity-pressure pattern remains large. The object to follow has two levels: one window \(Q_r\) at one scale, and the same solution seen through the nested windows \(Q_r\) as \(r\) runs down to zero.

\section{From One Cylinder To A Shrinking Family}

For one chosen radius \(r\), the region \(Q_r\) is a genuine parabolic cylinder:
\[
Q_r=B_r(x_0)\times(T-r^2,T).
\]
Every time slice in that one region has the same spatial cross-section:
\[
S(t)=B_r(x_0).
\]
Stacking \(B_r,B_r,B_r,\ldots\) through the interval \((T-r^2,T)\) gives one cylinder.

The singularity test changes the radius. A larger cylinder gives a coarse measurement; a smaller cylinder gives a sharper one; the limiting question follows the same solution through all of those measurements. In that cross-scale view the spatial slice is a ball whose radius changes with time:
\[
S(t)=B_{\rho(t)}(x_0),
\qquad
\rho(t)\to0\quad(t\to T).
\]
Stacking those time-dependent slices gives a funnel.

For Navier-Stokes the natural shrinking law is parabolic:
\[
\rho(t)=\lambda\sqrt{T-t}.
\]
The corresponding spacetime object is
\[
\F=
\left\{(x,t):0<T-t<\tau_0,\ |x-x_0|<\lambda\sqrt{T-t}\right\}.
\]
Its wall is
\[
\Sig=
\left\{(x,t): |x-x_0|=\lambda\sqrt{T-t}\right\}.
\]

The funnel is the spacetime track of the shrinking balls. It keeps the suspected point fixed, lets the observation scale collapse, and follows the last parabolic amount of time before \(T\). The question is whether the same fluid can carry one coherent scale-critical profile through that shrinking track.

\section{The Funnel Turns Finite-Time Blow-Up Into A Profile Problem}

Set
\[
\tau=T-t,
\qquad
\rho(t)=\sqrt{\tau},
\qquad
y=\frac{x-x_0}{\rho(t)},
\qquad
s=-\log\tau .
\]
Define the rescaled fields
\[
v(y,s)=\rho(t)u(x,t),
\qquad
q(y,s)=\rho(t)^2p(x,t).
\]
Then the funnel becomes the fixed cylinder in rescaled variables:
\[
|y|<\lambda,
\qquad
s\to\infty.
\]
A physical velocity of size
\[
u\sim (T-t)^{-1/2}
\]
becomes
\[
v\sim O(1).
\]
So the alleged terminal singularity becomes a persistent profile \(v(y,s)\) evolving for infinite rescaled time.

In these variables the Navier-Stokes equation has the schematic form
\[
\partial_s v-\nu\Delta_y v+(v\cdot\nabla_y)v+\nabla_y q
+\frac12 y\cdot\nabla_y v+\frac12 v=0,
\qquad
\nabla_y\cdot v=0.
\]
The terms
\[
\frac12 y\cdot\nabla_y v,
\qquad
\frac12 v
\]
are the cost of using the shrinking frame.

In this frame, terminal concentration becomes a question about the long-time behavior of \(v\) on \(B_\lambda\times(S_0,\infty)\). Regularity follows only if the rescaled profile has a scale-critical budget that eventually makes the late slabs small.

\section{The Moving Wall Is A Symptom, Not The Whole Problem}

Let
\[
e=\frac12|u|^2,
\qquad
\Omega(t)=B_{\rho(t)}(x_0).
\]
For a smooth pre-terminal solution, the local energy equation is
\[
\partial_t e+\nabla\cdot((e+p)u)=\nu\Delta e-\nu|\nabla u|^2.
\]
Reynolds transport on the moving ball gives the schematic identity
\[
\frac{d}{dt}\int_{\Omega(t)}e\dd x
+\nu\int_{\Omega(t)}|\nabla u|^2\dd x
=
\int_{\partial\Omega(t)}
\left[
 e\rho'(t)
 -(e+p)u\cdot n
 +\nu\partial_n e
\right]\dd S .
\]
The four wall channels are moving-wall loss, transport flux, pressure work, and viscous wall exchange.

For the parabolic funnel,
\[
\rho(t)=\lambda\sqrt{T-t},
\qquad
\rho'(t)=-\frac{\lambda}{2\sqrt{T-t}}.
\]
The wall speed diverges near the tip. The moving boundary changes the accounting, while the cost still has to be measured with the shrinking surface area included.

Write \(\tau=T-t\). Under the critical scaling
\[
\rho\sim\sqrt{\tau},
\qquad
|u|\sim\rho^{-1}\sim\tau^{-1/2},
\]
we have
\[
e\sim\tau^{-1},
\qquad
|\rho'|\sim\tau^{-1/2},
\qquad
|\partial B_\rho|\sim\rho^2\sim\tau.
\]
The moving-wall term per unit physical time scales like
\[
e|\rho'|\,|\partial B_\rho|
\sim
\tau^{-1}\tau^{-1/2}\tau
=
\tau^{-1/2}.
\]
The integral
\[
\int_0^{\tau_0}\tau^{-1/2}\dd\tau
\]
is finite.

The naive wall argument fails here. The physical wall payment may stay finite even as the wall accelerates. The remaining danger is scale-critical: a persistent rescaled profile can keep paying on every \(s\)-slab while the physical energy seen in shrinking variables is discounted.

\section{The Two Currencies}

In rescaled variables, the physical energy in the funnel is
\[
E_{\mathrm{phys}}(t)
=
\int_{B_{\lambda\sqrt{\tau}}}|u|^2\dd x
=
\tau^{1/2}\int_{B_\lambda}|v|^2\dd y
=
e^{-s/2}\int_{B_\lambda}|v|^2\dd y.
\]
A persistent \(O(1)\) rescaled profile can have physical energy tending to zero.

Physical dissipation carries the same kind of discount:
\[
\int |\nabla_xu|^2\dd x\dd t
=
\int e^{-s/2}|\nabla_yv|^2\dd y\dd s.
\]
The CKN-relevant cone quantity removes that discount:
\[
C_\lambda([S,S+1])
=
\int_S^{S+1}\int_{B_\lambda}
\left(|v|^3+|q|^{3/2}\right)\dd y\dd s.
\]

The forward theorem would need to prove an undiscounted budget such as
\[
\int_{S_0}^{\infty}\int_{B_\lambda}
\left(|v|^3+|q|^{3/2}+|\nabla_yv|^2\right)\dd y\dd s
<\infty.
\]
Then the late slabs satisfy
\[
\int_S^{S+1}\int_{B_\lambda}
\left(|v|^3+|q|^{3/2}\right)\dd y\dd s
\to0
\qquad(S\to\infty),
\]
and CKN regularity controls the tip.

The hard bridge is discount removal:
\[
\text{finite discounted physical funnel accounting}
\quad\Longrightarrow\quad
\text{undiscounted scale-critical cone control}.
\]
Pure scaling does not give this. It is exactly where Navier-Stokes structure must enter.

\section{The Correct Order Of The Proof Program}

The funnel should first be treated as a forward-positive control object. The order is:
\[
\begin{aligned}
&\text{funnel local energy inequality}\\
&\quad\to \text{discounted cone budget}\\
&\quad\to \text{attempt discount removal}\\
&\quad\to \text{precise surviving obstruction}.
\end{aligned}
\]
Only after that obstruction is known should it be classified as a CM class-exit witness.

This matters because a premature class-exit label can hide the actual estimate that has to be tried. The forward children are:
\[
\boxed{ParabolicFunnelScaleCriticalBudget.A}
\]
and
\[
\boxed{FunnelProfileLyapunov.A}.
\]
The first would prove the undiscounted scale-critical cone budget directly. The second would find a monotone or coercive functional for the rescaled profile equation that forces late-slab smallness.

If those forward attempts fail in a precise way, the surviving obstruction can be sent downstream to the Pack/Part/Field witness tree. The classification receives an obstruction; it does not replace the forward analysis.

\section{The Density-Only Objection}

Region mass vanishes with volume, so region-mass language obscures the pointwise velocity question. The statement should be written in densities.

Use mass-density language directly:
\[
\rho(x,t)=\rho_0>0.
\]
Define momentum density and kinetic-density by
\[
\pi(x,t)=\rho_0u(x,t),
\qquad
k(x,t)=\frac12\rho_0|u(x,t)|^2.
\]
The pointwise velocity relation is exact:
\[
|u(x,t)|^2=\frac{2k(x,t)}{\rho_0}.
\]
So
\[
0<\rho_0<\infty
\quad\text{and}\quad
k(x_0,T)<\infty
\quad\Longrightarrow\quad
|u(x_0,T)|<\infty.
\]
This algebraic step is closed.

The open problem is the premise:
\[
k(x_0,T)<\infty.
\]
The standard energy control gives a regional integral:
\[
\int_{B_R(x_0)}k(x,T)\dd x<\infty.
\]
That does not automatically give the pointwise statement at \(x_0\).

\section{The Measurement-Order Correction}

The transcript's sharpest density correction was about measurement order. The wrong shortcut is to shrink the measuring ball first and then read the resulting small integral as control. The useful thought is different.

First take a fixed ball:
\[
E_R=\int_{B_R(x_0)}k(x,T)\dd x<\infty.
\]
Now ask what happens if that same finite amount \(E_R\) is collapsed into a smaller ball \(B_\varepsilon(x_0)\). The average kinetic-density becomes
\[
\overline{k}_\varepsilon
=
\frac{E_R}{|B_\varepsilon|}
\sim
\frac{E_R}{\varepsilon^3}.
\]
Thus
\[
\overline{k}_\varepsilon\to\infty
\qquad(\varepsilon\to0),
\]
and
\[
|u_\varepsilon|^2
\sim
\frac{2E_R}{\rho_0\varepsilon^3}
\to\infty.
\]

In the limit this is no longer an ordinary bounded density field. It becomes a measure-like concentration:
\[
k_\varepsilon(x)\dd x \to E_R\delta_{x_0}.
\]
Classical Navier-Stokes wants ordinary fields, not a hidden delta of kinetic-density at the endpoint. The proof must either rule out that concentration by participation estimates or classify it as a failure of the ordinary same-fluid field object.

This also explains why finite integrated energy is too weak by itself. A model density
\[
k(x)\sim |x-x_0|^{-\alpha},
\qquad
0<\alpha<3,
\]
has
\[
\int_{B_R(x_0)}k(x)\dd x<\infty
\]
while its point value at \(x_0\) is unbounded.

\section{The Theorem The Sample Wants}

The clean density theorem is a density-to-density upgrade. In funnel variables define
\[
K(y,s)=\frac12\rho_0|v(y,s)|^2.
\]
A useful closure theorem would say:
\[
\boxed{
\text{finite same-fluid scale-critical density ledger on }\F
\Longrightarrow
\sup_{\substack{s\ge S_0\\ |y|<\lambda/2}}K(y,s)<\infty .
}
\]
Then
\[
|v(y,s)|^2=\frac{2K(y,s)}{\rho_0}
\]
is bounded in the late funnel core, and the terminal velocity is controlled after translating back to physical variables.

A Morrey-style sufficient form is
\[
\boxed{KineticDensityMorreyFromParticipation.A}
\]
with
\[
\int_{B_r(x_0)}k(x,T)\dd x
\le
C|B_r|
\qquad
\text{for all small }r.
\]
Equivalently,
\[
\sup_{0<r<r_0}
\frac1{|B_r|}
\int_{B_r(x_0)}k(x,T)\dd x
<\infty.
\]
That is the exact upgrade from finite regional energy to finite pointwise kinetic-density scale.

The matching funnel theorem is
\[
\boxed{ParabolicFunnelScaleCriticalBudget.A}
\]
with estimate
\[
\int_{S_0}^{\infty}\int_{B_\lambda}
\left(|v|^3+|q|^{3/2}+|\nabla_yv|^2\right)\dd y\dd s
<\infty.
\]
If this holds, late cone slabs calm down:
\[
\int_S^{S+1}\int_{B_\lambda}
\left(|v|^3+|q|^{3/2}\right)\dd y\dd s
\to0,
\]
and the CKN criterion can be applied near the tip.

\section{Conclusion}

The fixed-cylinder picture fails because it keeps the radius constant while the singularity question sends the radius to zero. The parabolic funnel is the spacetime region traced by that collapse. In the funnel variables, terminal concentration becomes a persistent profile problem.

The naive wall argument also fails. A fast-moving wall does not automatically create infinite physical cost, because the wall area is shrinking. The obstruction has to be stated in scale-critical currency: a persistent profile may keep a non-small CKN packet on every late slab while the physical energy is discounted.

The density endpoint is simpler. Positive mass-density does not by itself bound velocity. It bounds velocity only after pointwise kinetic-density is finite:
\[
k(x_0,T)<\infty\Longrightarrow |u(x_0,T)|<\infty.
\]

Finite regional kinetic energy is weaker than that pointwise premise. Collapsing a fixed finite energy into zero volume creates a singular kinetic-density concentration, not an ordinary bounded field value. The proof must obtain a scale-uniform density bound, such as a Morrey-type estimate or a finite undiscounted funnel budget.

The resulting chain is:
\[
\begin{aligned}
&\text{shrinking CKN packet}\\
&\quad\to \text{parabolic funnel}\\
&\quad\to \text{rescaled persistent profile}\\
&\quad\to \text{discounted versus scale-critical budget}\\
&\quad\to \text{density-to-density point closure}.
\end{aligned}
\]

\end{document}
